\section{Conclusion}
\label{sec:conclusion}

We presented a Systematization of Knowledge for metadata-hiding group
registries with SNARK-gated state changes, structured as a six-axis
taxonomy that indexes by deployment constraint rather than by
construction family. Across seven concrete configurations and five
comparison dimensions, the dominant finding is that no configuration
sits in every priority's preferred regime: the cheap-and-not-PQ
regime (C1, C2), the PQ-ready-but-large-proof regime (C6), and the
host-function-blocked hypothetical (C7) form a partial order, not a
total one. The deployment recommendations of
\S\ref{sec:recommendations} are accordingly profile-conditional, and
two of the four profiles point to configurations that would not be
selected under the profile-blind reading of the comparison tables.

The contribution we believe is the most actionable for protocol
designers is the cross-axis dependency analysis of
\S\ref{sec:dimensions-deps}: fixing the anchor chain for
non-cryptographic reasons implicitly fixes four of the six axes
through the chain~$\to$~host~$\to$~curve~$\to$~SNARK~$\to$~setup
chain, leaving only the in-circuit hash truly free. Designers who
treat the chain choice as independent of the proof-system choice
underestimate the constraint they have already accepted.

The contribution we believe is the most actionable for the
cryptographic-research community is the migration matrix of
\S\ref{sec:migration-matrix} and the open-problem cluster of
\S\ref{sec:open}: no column in the migration matrix is uniformly
Class~(a), and the empty cheap-and-PQ-ready cell of
\S\ref{sec:analysis-synthesis} is closed only by progress on either
chain-host-function ABIs that admit FRI primitives, or on
SNARK constructions that combine pairing-friendly verifier costs
with hash-based post-quantum security. The path of least research
risk is the chain-host axis; the path of greatest impact is the
SNARK-construction axis.

Future work in two directions is concrete enough to commit to a
follow-up SoK. First, an updated taxonomy that admits a fourth
post-quantum class --- lattice-based --- once a deployable
construction in that class exists. Second, an extension of the
abstraction to operator-aware metadata-hiding against
jurisdiction-bound adversaries (open problem~8); the threat model in
\S\ref{sec:problem} explicitly excludes this adversary class today,
and the taxonomy will need a corresponding extension before that
extension can be made.
